\documentclass[border=6pt]{standalone}
\usepackage{amsmath}
\usepackage{tikz}
\usetikzlibrary{calc,arrows.meta}

\begin{document}
\begin{tikzpicture}[scale=1, >=Stealth]
  \draw[help lines, color=gray!30] (-1,-1) grid (6,4);
  \draw[->] (-1,0) -- (6,0) node[right] {$x$};
  \draw[->] (0,-1) -- (0,4) node[above] {$y$};

  % Segitiga asal (biru)
  \coordinate (P) at (1,1);
  \coordinate (Q) at (4,1);
  \coordinate (R) at (2,3);

  % Segitiga bayangan (dihitung dari vektor pergeseran (3,2))
  \foreach \p in {P,Q,R} {
    \coordinate (\p') at ($(\p)+(3,2)$);
  }

  % Vektor pergeseran dari tiap titik sudut
  \foreach \p in {P,Q,R} {
    \draw[very thick, -Stealth, yellow!85!black] (\p) -- (\p');
  }

  \draw[blue, very thick] (P) -- (Q) -- (R) -- cycle;
  \draw[orange, very thick] (P') -- (Q') -- (R') -- cycle;

  \foreach \p/\l/\pos in {P/{P(1,1)}/below left, Q/{Q(4,1)}/below right, R/{R(2,3)}/above left} {
    \filldraw[blue] (\p) circle (2pt) node[\pos] {\l};
  }
  \foreach \p/\l/\pos in {P'/{P'(4,3)}/left, Q'/{Q'(7,3)}/below right, R'/{R'(5,5)}/above right} {
    \filldraw[orange] (\p') circle (2pt) node[\pos] {\l};
  }
\end{tikzpicture}
\end{document}
